% Fichier complété par tools/complete_missing_corrections.py.
% Source : SYS/SYS-01/SYS-01_ChaineInfo/538_Codeur/538_Codeur.tex
% Statut : rédigé (5 question(s)).
\begin{corrigebox}[Corrigé rédigé]
\CorrigeQuestion{1}
Un codeur incrémental optique comporte un disque à fentes, une source et
            un photodétecteur. La rotation produit des impulsions; deux voies A/B en
            quadrature donnent le sens et une voie index fournit une référence par tour.
\CorrigeQuestion{2}
Avec 25 fentes et une seule voie, la résolution est
            $360^\circ/25=\boxed{14,4^\circ}$ par impulsion. En comptage des deux fronts,
            elle devient $7,2^\circ$, et avec deux voies/quatre fronts $3,6^\circ$.
\CorrigeQuestion{3}
La fréquence des impulsions vaut $f=N n/60$ pour $N$ impulsions par tour
            et une vitesse $n$ en tr/min. La vitesse est donc $n=60f/N$; le sens provient
            de l'ordre des fronts A et B.
\CorrigeQuestion{4}
La position relative est le nombre algébrique d'impulsions multiplié par
            la résolution. Une prise d'origine sur la voie index est nécessaire après
            mise sous tension.
\CorrigeQuestion{5}
L'incertitude de quantification est au minimum d'un demi-pas, soit
            $\pm7,2^\circ$ avec une seule voie/25 fentes, avant prise en compte des erreurs
            optiques et mécaniques.
\end{corrigebox}
